\documentclass[11pt]{article}
\usepackage[margin=1in]{geometry}
\usepackage{amsmath,amssymb,amsthm,booktabs}
\usepackage[hidelinks]{hyperref}
\setlength{\emergencystretch}{2em}
\newtheorem{theorem}{Theorem}
\newtheorem{proposition}[theorem]{Proposition}
\newtheorem{corollary}[theorem]{Corollary}
\newcommand{\mr}{\operatorname{mr}}
\newcommand{\rk}{\operatorname{rank}}
\newcommand{\Hom}{\operatorname{Hom}}
\newcommand{\Z}{\mathbb Z}
\newcommand{\F}{\mathbb F}
\title{Meridional lift bounds through branched covers}
\author{Internal research working proof}
\date{16 November 2026}
\begin{document}
\maketitle

For a transposition quotient $\phi:G_K\twoheadrightarrow S_r$, the
restricted affine lift space has a topological interpretation:
\[
 \dim_F Z_\phi(F)=r-1+\dim_F H_1(Y_\phi;F),
\]
where $Y_\phi$ is its connected simple $r$-sheeted branched cover of
$S^3$. A graph argument strengthens the resulting homological bound to
\[
                 \mr(K)\ge r-1+\rk\pi_1(Y_\phi).
\]
Here rank of a group means its minimum number of arbitrary generators.
These are written proofs, with finite controls described below; they
have not received formal verification or external specialist review.
No priority claim is made. In degree two the nonabelian construction is
the classical Kelly--Wada core-group method \cite{core}.

The new interpretation also closes a route on the retained sixteen
thirteen-crossing targets: no symmetric root lift over any field or
degree can be sharp there. The original binary common-form criterion
is impossible on these targets in every dimension. Neither conclusion
resolves their MRC status. The full conjecture remains open.

\section{The branched-cover group bound}

Fix $r\ge2$. Complete the natural connected $r$-sheeted cover of the
knot exterior defined by $\phi$ by filling every lifted meridian cycle.
A meridian has one cycle of length two and $r-2$ fixed sheets. The
corresponding filling slopes are its square on the first component and
its first power on each fixed-sheet component. The resulting closed
oriented manifold is $Y_\phi$.

\begin{theorem}\label{thm:group}
For every such quotient,
\[
                     \mr(K)\ge r-1+\rk\pi_1(Y_\phi).
\]
\end{theorem}
\begin{proof}
Choose any $w$ meridional generators of $G_K$, and a presentation on
these generators. Its lifted presentation complex has $r$ vertices
and, for each generator with image $(pq)$, $r$ directed edges. On each
fixed sheet the edge is a loop killed by the meridian filling. The two
remaining edges run from $p$ to $q$ and back. Filling the lifted square
identifies these two edges as inverse paths. This leaves a graph with
$r$ vertices and precisely $w$ unoriented edges, one per generator.
The graph is connected because these transpositions generate $S_r$.

A maximal tree uses $r-1$ edges. Collapsing it leaves $w-r+1$ potential
generators. Lifted presentation relators and all other filling relations
can only impose relations, so $\pi_1(Y_\phi)$ is a quotient of a free
group on these $w-r+1$ generators. Minimizing over $w$ proves the claim.
The argument is unchanged if some chosen meridians have inverse
orientation: their projected transpositions and filled cycles are the
same. A presentation complex suffices for this computation of the
fundamental group; no asphericity assertion is needed.
\end{proof}

\section{What the integer lift matrix computes}

Let $m_1,\ldots,m_s$ be meridional presentation generators and
$\phi(m_i)=(p_iq_i)$, with $p_i<q_i$. For a field $F$ put
$E_r(F)=\{x\in F^r:\sum_jx_j=0\}$. The companion all-characteristic
note defines $Z_\phi(F)$ by the lifts
\[
 m_i\longmapsto (a_i(e_{p_i}-e_{q_i}),(p_iq_i))
                       \quad\hbox{in }E_r(F)\rtimes S_r.
\]
The translation coordinates of all relators give an integer matrix
$M_\phi$ with $s$ columns, and $Z_\phi(F)=\ker(M_\phi\otimes F)$.
The matrix need not have a square or minimal form.

\begin{theorem}\label{thm:homology}
There is an exact sequence of vector spaces
\[
 0\longrightarrow F^r/F\mathbf1
   \longrightarrow Z_\phi(F)
   \longrightarrow H^1(Y_\phi;F)\longrightarrow0.
\]
Thus $\dim Z_\phi(F)=r-1+\dim H_1(Y_\phi;F)$ in every characteristic,
including those dividing $r$. The nonunit Smith invariant factors of
$M_\phi$ are exactly the torsion invariant factors of
$H_1(Y_\phi;\Z)$. If $t=\rk_\Z M_\phi$, then
\[
       \rk_\Z H_1(Y_\phi;\Z)=s-r+1-t.
\]
\end{theorem}
\begin{proof}
Use the lifted edges and fillings in Theorem~\ref{thm:group}. After
eliminating the fixed loops and one edge from each moved pair, their
integral chain group in degree one is $\Z^s$. Its boundary sends the
$i$th oriented edge to $e_{q_i}-e_{p_i}$. Write this connected
multigraph as $\Gamma$ and its cycle lattice as
$L=\ker(\partial:\Z^s\to\Z^r)$.

A scalar $a_i$ assigns opposite translations to the two orientations
of this edge, and zero to the killed fixed-sheet loops. These are
precisely the displayed root translations, up to the harmless choice
of edge orientation. Evaluating a lifted relator adds its edge values
along the lifted path. Its vanishing on every starting sheet is
exactly the corresponding affine relator equation. Path conventions
may reverse a relator or an edge; these change neither relation lattice
nor kernel. Thus the rows of $M_\phi$, up to these choices, generate
the subgroup $R\subset L$ of lifted relator chains.

Over $F$, the solutions are the cellular one-cocycles of this filled
complex. The coboundary of a vertex potential $b\in F^r$ assigns the
edge difference $b_{q_i}-b_{p_i}$. Connectivity makes its kernel
exactly $F\mathbf1$, yielding the exact sequence. These are cuts from
\emph{all} vertex potentials. Restricting potentials to $E_r(F)$ would
lose cuts when the characteristic divides $r$ and would give a wrong
formula.

For the integral assertion, choose a spanning tree in $\Gamma$.
Projection from $L$ to the $s-r+1$ chord coordinates is an integral
isomorphism: a chord assignment has a unique integral completion to a
cycle, obtained along the tree. Also $L$ is a direct summand of
$\Z^s$, since $\Z^s/L\cong\operatorname{im}\partial$ is free. Hence an
integral change of basis puts the row presentation into $r-1$ zero
columns followed by a presentation of $L/R=H_1(Y_\phi;\Z)$. This proves
both Smith and free-rank assertions.

The fillings are essential: merely abelianizing the unfilled exterior
cover computes a different group. Since every displayed meridian lift
is an involution, replacing inverse meridian letters by positive ones
in $M_\phi$ is valid. The signed lifted paths in the exterior become
equivalent only after the stated meridian-cycle fillings.
\end{proof}

In particular, the all-characteristic maximum equals
\[
 r-1+\rk_\Z H_1(Y_\phi;\Z)
       +\#\{\hbox{nonunit torsion invariant factors of }H_1(Y_\phi;\Z)\}.
\]
A prime divisor of the first nonunit factor attains it, by divisibility
of the invariant factors. Linear lifts see the abelianization of the
cover group, while Theorem~\ref{thm:group} retains its nonabelian rank.

\section{A nonabelian lift count}

Let $A$ be a finite group with $|A|>1$. In the wreath product
$A^r\rtimes S_r$, restrict a meridian over $(pq)$ to have coordinates
$a$ at $p$, $a^{-1}$ at $q$ and identity elsewhere. Its square is the
identity. These elements form a conjugation-invariant set, with
exactly $|A|$ choices over each transposition. Inversion preserves it.
Let $\mathcal L_\phi(A)$ denote the homomorphisms lifting the fixed
$\phi$ subject to this restriction.

\begin{theorem}\label{thm:nonabelian}
\[
 |\mathcal L_\phi(A)|=|A|^{r-1}|\Hom(\pi_1(Y_\phi),A)|,
 \qquad
 \mr(K)\ge r-1+
 \left\lceil\log_{|A|}|\Hom(\pi_1(Y_\phi),A)|\right\rceil.
\]
\end{theorem}
\begin{proof}
The coordinates are group-valued transports along the edges of
$\Gamma$, with inverse transports along reverse edges. The wreath
relator equations say exactly that every lifted relator has trivial
transport. The fixed-sheet and square conditions are the meridian
fillings. Thus these lifts are flat $A$-valued edge assignments on the
filled complex.

Fix a root vertex and a spanning tree. Vertex changes of trivialization
with identity at the root form a group of order $|A|^{r-1}$. They act
freely: preservation of the edge assignment propagates the identity
from the root along the tree. Every assignment has a unique such
change making every tree transport the identity. The remaining chord
transports then give exactly a homomorphism $\pi_1(Y_\phi)\to A$,
without quotienting by conjugacy. This proves the count.

Evaluation on any $w$ meridional generators injects
$\mathcal L_\phi(A)$ into a set of size $|A|^w$. Taking logarithms and
minimizing $w$ proves the bound. It also follows from
Theorem~\ref{thm:group}, since a $d$-generated group has at most $|A|^d$
homomorphisms to $A$.
\end{proof}

For $r=2$ every knot has the cyclic transposition quotient, and
$Y_\phi=\Sigma_2(K)$. A crossing with over-colour $b$ and incoming
under-colour $a$ has outgoing colour $ba^{-1}b$. These are the
nonabelian core colourings. The core group is
$\pi_1(\Sigma_2(K))*\Z$; fixing any one arc colour to the identity
counts $\Hom(\pi_1(\Sigma_2(K)),A)$ \cite{core}. This historical result
is not a new theorem of the present work.

\section{A geometric sharpness gate}

A $b$-bridge decomposition lifts to a Heegaard decomposition of
$Y_\phi$ of genus $b-r+1$. Each side is the cover of a trivial
$b$-string tangle. It is connected since its bridge meridians generate
$G_K$ and hence $S_r$; the cover is a handlebody obtained from $r$
balls and $b$ one-handles. Consequently
\[
       \rk\pi_1(Y_\phi)\le g(Y_\phi)\le b(K)-r+1.
\]
Here $g$ denotes minimum Heegaard genus. A certified lower bound
$\rk\pi_1(Y_\phi)\ge b-r+1$ therefore proves MRC when accompanied by
a $b$-bridge upper certificate. A sufficient homological condition is
$\dim_FH_1(Y_\phi;F)=b-r+1$.

Neither equality in this chain is established for arbitrary knot
covers. Turning this gate into a universal existence theorem would
require additional mathematics. The original affine criterion already
fails on the infinite torus family proved in the companion note.

For any subgroup $J\le\pi_1(Y_\phi)$ of finite index $d$, there is also
the bound
\[
                 \mr(K)\ge r+\left\lceil\frac{\dim_F H_1(J;F)-1}{d}\right\rceil.
\]
Indeed, a $t$-generator presentation for $\pi_1(Y_\phi)$ lifts to a
connected $d$-vertex graph with $dt$ edges. Collapsing a spanning tree
gives at most $d(t-1)+1$ generators for $J$, the usual Schreier bound.
Hence $\dim H_1(J;F)\le d(t-1)+1$, and Theorem~\ref{thm:group}
gives the formula. This does not assume that $J$ is normal or that
homology detects its rank. For $r=2,d=5$, dimension at least seven
suffices for $\mr\ge4$.

\section{Complete exclusions on the retained current panel}

The same sixteen targets were frozen from the first sixteen names in
the pinned October 2026 uncertainty list \cite{deposit}. Their oriented
PD data come from the independent KnotInfo 2026.11.2 archive. Every
diagram has a retained four-seed bridge certificate. The results below
concern these inputs; the uncertainty-list label is not a novelty proof.

\begin{proposition}\label{prop:all}
For these sixteen diagrams no symmetric transposition root lift over
any field or any degree gives a lower bound exceeding two. No
transitive binary common-form affine representation satisfying the
original quadratic criterion exists in any dimension.
\end{proposition}
\begin{proof}
If $w\le4$ meridians generate $G_K$, their $S_r$ transpositions form
a connected graph on $r$ vertices with $w$ edges. Therefore $r\le5$.
It suffices to test $r=2,3,4,5$, without an arbitrary degree cutoff.
The retained exhaustive normalized four-seed searches find one $S_2$
colouring and no surjective $S_3,S_4,S_5$ colouring in every case.
An independent permutation enumeration and propagation replay checks
this conclusion. Each $S_2$ root matrix has thirteen columns and
nonzero Smith factors $1^{11},D$, with $D>1$ below. Its maximum
nullity over all fields is therefore exactly two. Higher-degree
transposition quotients cannot exist, and no choice of characteristic
can improve this method's bound on the stated panel.

For the original binary criterion, write its vector-space dimension
as $k$. Its proved lower bound and the four-seed upper bound give
$k+1\le\mr(K)\le b(K)\le4$, so $k\le3$. A nonzero alternating form
in those dimensions necessarily has rank two. The unlimited-radical
proposition in the companion all-characteristic note shows that such
a transitive action forces an $S_3$ transposition quotient, absent
here. Thus the original route is excluded in every dimension. The
earlier suggestion that higher-rank binary forms might work on these
sixteen targets is incorrect and is withdrawn.
\end{proof}

\begin{center}
\begin{tabular}{rrrrrr}
\toprule
Knot & $D$ & $|\Hom(\pi_1\Sigma_2,A_5)|$ &
Knot & $D$ & $|\Hom(\pi_1\Sigma_2,A_5)|$\\
\midrule
$13a129$&391&121&$13a233$&403&1\\
$13a135$&401&241&$13a234$&389&121\\
$13a136$&391&121&$13a239$&361&1\\
$13a137$&401&241&$13a242$&359&121\\
$13a194$&391&121&$13a243$&337&1\\
$13a196$&401&241&$13a244$&311&121\\
$13a197$&391&121&$13a253$&337&1\\
$13a198$&401&241&$13a265$&311&121\\
\bottomrule
\end{tabular}
\end{center}

In each case $H_1(\Sigma_2;\Z)=\Z/D$. The $A_5$ counts use a separately
frozen search: normalize the first seed to the identity and enumerate
all $60^3=216000$ assignments on the remaining three seeds, per case.
To prove $\mr\ge4$ by this count requires strictly more than
$60^2=3600$ normalized colourings. The observed counts $1,121,241$
do not do so. A separate $\operatorname{PSL}_2(\F_5)$ matrix
implementation reproduces all sixteen counts. These finite results
exclude this particular counting certificate, not other coefficient
groups or the stronger nonabelian cover-rank gate.

The user's subsequent joint-kernel test strengthens this approach:
twenty distinct $A_5$ quotient kernels of one double-cover group
force its rank to be at least three, hence $\mr\ge4$. The classical
capacity is nineteen generating-pair automorphism classes, stated
explicitly by Burness--Crestani with credit to Hall \cite[p.~2]{bc}.
The companion \path{JOINT_SIMPLE_QUOTIENT_KERNELS.tex} proves the
general capacity bound and separately replays all 293 retained
uncertain names: the maximum is five kernels on one knot, so none
is sharp. The twelve first witnesses above cannot be combined across
knots. A known positive control, $P(2,3,3,5)$, has 23 kernels and
2901 homomorphisms. The kernel test gives sharp four there, while
the raw count gives only three. It is already a settled Montesinos
case, not a new resolution.

A second frozen test uses exactly the first nonconstant $A_5$ witness
already retained for each target, without replacing witnesses after
failure. Twelve targets have such a witness, with image all of $A_5$;
the other four have no nonconstant map by the complete count above.
The natural action on five letters supplies an index-five subgroup
$J\le\pi_1(\Sigma_2K)$. Its presentation is computed from the core
presentation after killing the normalized first arc. All twelve
degree-five covers have $H_1(J;\Z)=\Z/(3D)$, with $D$ in the table.
Their maximum homology dimension over every field is only one,
so none reaches the Schreier threshold seven. Only these retained
first witnesses are excluded; other covers remain untested.

The user's stronger augmentation-module test now completes this
degree-five boundary. For an index-$n$ subgroup $J\le\pi_1Y$ and a
field with characteristic not dividing $n$, Shapiro and the splitting
of the permutation module give
\[
 \mr(K)\ge r+\left\lceil
 \frac{b_1(J;F)-b_1(\pi_1Y;F)}{n-1}\right\rceil.
\]
For $r=2,n=5$, gain at least five suffices for sharp four. The
companion \path{AUGMENTATION_COVER_HOMOLOGY.tex} tests all sixteen
$A_5$ kernel classes on these same sixteen knots, including four
classes beyond the twelve earlier witnesses. Every subgroup has
$H_1=\Z/(3D)$, giving gain one at three and zero at every other
prime. Since every $D$ is coprime to thirty, any transitive
$G\to S_5$ has image exactly $A_5$: sign is trivial and proper
$A_5$ images would force an absent abelian quotient at two, three
or five. Thus every connected index-five cover is represented.
The proposed Betti-gain certificate fails in degree five on this
panel; other degrees and nonabelian subgroup information remain open.

These same Smith certificates also exclude the explicit
Lustig--Moriah gcd plat criterion in bridge index $n\ge3$ on this
panel. To see this directly, propagate Fox colours through a
$2n$-plat: the unreduced Burau matrix at $t=-1$ gives the strand
colours, $M$ duplicates the $n$ upper-pair colours, and $L$ tests
equality at the lower pairs, with $L,M$ as in the literature note
\cite{moriah}. Thus the Fox-colouring space is
$\ker(L\widehat\rho(B)M)$ over its coefficient field. If a prime
divides every matrix entry, this space has dimension $n$. That is
the degree-two root-lift space, whose all-field dimension is at most
two here. A nonunit gcd therefore cannot occur for $n\ge3$.
This is a consequence of the linear obstruction, not an audit of
every geometric generalized-Montesinos convention.

\section{Verification, provenance and limits}

\texttt{verify\_branched\_cover\_bounds.py} independently enumerates
the projected groups and checks all crossings after seed propagation.
It replays the $S_2$ Smith operations against direct Fox rows
$2a_{\rm over}-a_{\rm in}-a_{\rm out}$, and checks dimensions in
characteristics $2,3,5,7,11$. The all-degree conclusion additionally
uses the proved $r\le5$ bound, and the all-field conclusion uses
integer Smith factors, not these five sample primes.

For homology controls, \texttt{branched\_cover\_homology.py} uses all
$rs$ directed presentation-cover edges and retains the signs of the
supplied relator words, then adds every fixed-sheet meridian and
moved-pair square filling. Diagram relators may first be rewritten
using the meridian-square relations, valid for this filled complex.
It does not construct the scalar root matrix. After collapsing
a spanning tree, the full-chain Smith computation agrees with the
root-matrix torsion and free-rank formula on nine controls:
\begin{itemize}
\item For the connected sum of $t=1,2,3,4$ trefoils, the double cover
has $H_1=(\Z/3)^t$ and the specified $S_3$ cover has
$H_1=\Z^{t-1}$.
\item For the retained sharp $S_4$ colouring of $15a1381$, the cover
has $H_1=\Z\oplus\Z/5$. Thus the root-lift dimensions are four in
characteristic zero and five in characteristic five, explaining the
earlier sharp certificate on this already-known Montesinos knot.
\end{itemize}
Every Smith computation is checked by replaying its elementary
unimodular operations. The two homology paths share the Smith producer
and replay routine; their agreement is a chain-construction control,
not independent verification of the uniform topological proof.

The nonabelian verifier instead constructs
$\operatorname{PSL}_2(\F_5)$ from $2\times2$ matrices modulo $\pm I$,
using its standard isomorphism with $A_5$. It checks order sixty and
element-order multiplicities $1,15,20,24$ for orders $1,2,3,5$.
The trefoil control gives twenty-one normalized core colourings,
as required by $\Hom(\Z/3,A_5)$. Its coefficient group and propagation
implementation are separate from the permutation producer; the PD
parser is shared. Receipts seal all inputs in
\texttt{branched\_cover\_verification.json} and
\texttt{nonabelian\_core\_verification.json}.

For the twelve index-five covers,
\texttt{verify\_core\_subgroup\_homology.py} checks the $A_5$ image,
every core relation and the normalization directly with permutations.
It constructs a separate integer chain matrix from evaluated Fox
prefix derivatives and replays every Smith step. The positive edge
action is the right monodromy $j\cdot g=\phi(g)^{-1}j$, so inverse
letters are retained. Matrices have seventy rows and sixty-one
columns after collapsing the four-edge tree. A free-group control
recovers rank six in an index-five subgroup of $F_2$; a Kurosh control
for $C_3*C_3$ gives $H_1=(\Z/3)^4$. An invalid relator is rejected.
The PD parser and Smith producer/replay are shared. The receipt
\texttt{core\_subgroup\_homology\_verification.json} seals the frozen
plan, the first-witness choice and all inputs. This finite negative
test is separate from the proof of the uniform Schreier inequality.

The exact known-family scopes added to the literature note close
specific audit gaps: the many-twig theorem requires at least three
added twigs at each vertex of the core tree and excludes weights
$0,\pm1$; the Lustig--Moriah plat criterion requires its integer gcd
to differ from one. Neither is an all-arborescent theorem. No new
previously unresolved knot is settled here. Further progress needs
a new lower-bound certificate or a uniform mechanism beyond the
excluded routes, together with the remaining literature and proof
audits.

\begin{thebibliography}{9}
\small
\setlength{\itemsep}{2pt}
\bibitem{bc} Burness--Crestani, On the generating graph of direct
powers of a simple group, J. Algebraic Combinatorics 38 (2013),
329--350, author manuscript p.~2.
\url{https://seis.bristol.ac.uk/~tb13602/docs/BC_1.pdf}.
\bibitem{moriah} Moriah, Heegaard splittings of knot exteriors,
Geometry \& Topology Monographs 12 (2007), 191--232, Theorem 4.20.
\url{https://doi.org/10.2140/gtm.2007.12.191}.
\bibitem{core} Silver--Traldi--Williams, Core groups, J. Knot Theory
Ramifications 33 (2024), 2450013, introduction and Section 7.
The core-group construction is credited there to Kelly (1991) and
Wada (1992). \url{https://arxiv.org/abs/2311.02048v2}.
\bibitem{deposit} Bhattacharjee, The bridge index and the nonorientable
four-genus of knots with at most thirteen crossings, unrefereed release
1.3.0 (2026). \url{https://doi.org/10.5281/zenodo.23561157}.
Official data: \url{https://github.com/anaxia-research/conjecture-to-choose}.
\end{thebibliography}
\end{document}
